% ৩.১১ বুলিয়ান উপপাদ্য
\section{Boolean Theorems}

\begin{itemize}
\item Identity, Null, Idempotent, Complement laws
\item Involution, Commutative, Associative, Distributive
\item Absorption, Duality principle
\end{itemize}

\begin{learningbox}
এই টপিক শেষে শিক্ষার্থী পারবে:
\begin{itemize}
\item Boolean algebra-এর মৌলিক উপপাদ্য লিখতে।
\item theorem ব্যবহার করে expression সরল করতে।
\item duality principle-এর ধারণা ব্যাখ্যা করতে।
\end{itemize}
\end{learningbox}

\begin{center}
\begin{tabular}{|p{0.24\textwidth}|p{0.30\textwidth}|p{0.30\textwidth}|}
\hline
\textbf{আইন} & \textbf{OR form} & \textbf{AND form} \\
\hline
Identity & $A+0=A$ & $A\cdot1=A$ \\
\hline
Null & $A+1=1$ & $A\cdot0=0$ \\
\hline
Idempotent & $A+A=A$ & $A\cdot A=A$ \\
\hline
Complement & $A+\overline{A}=1$ & $A\cdot\overline{A}=0$ \\
\hline
Involution & $\overline{\overline{A}}=A$ & $\overline{\overline{A}}=A$ \\
\hline
\end{tabular}
\end{center}

\begin{examplebox}{সরলীকরণ}
$A+AB=A(1+B)=A\cdot1=A$। এটি absorption law-এর ব্যবহার।
\end{examplebox}

\begin{keypointbox}{Duality principle}
Boolean expression-এ $+$ ও $\cdot$ এবং 0 ও 1 পরস্পর বদল করলে যে dual expression পাওয়া যায়, সেটিও বৈধ theorem হতে পারে।
\end{keypointbox}
